% !PRINTED 471
\weberpage{471}

%<<CONT
was nur möglich ist, wenn $B \equiv 0\ (\mathrm{mod}\ 4)$ ist. Um aber zu entscheiden, welche der beiden Gleichungen (15) erfüllt ist, kommt es nach § 40, (12) darauf an, ob
\[
(\alpha - \beta)^2 - 1 + (\alpha - \beta)(\alpha + \beta + \gamma - \delta)
\]
oder, was dasselbe ist,
\[
(B+y)^2 - 1 + (B+y)(A - nC)
\]
durch 16 oder nur durch 8 teilbar ist. Vermehrt man aber in diesem Ausdruck $B$ um ein ungerades Vielfache von 4, so verändert er sich um ein ungerades Vielfache von 8, woraus zu schließen, daß, je nachdem $B \equiv 0$ oder $\equiv 4\ (\mathrm{mod}\ 8)$, die \gesperrt{eine oder die andere} der beiden Gleichungen (15) befriedigt ist. % [sic] printed phrase omits the grammatically required "ist"

Damit ist bewiesen:

6. \gesperrt{Ist $m \equiv 1\ (\mathrm{mod}\ 8)$, so ist $\sqrt{2}\, f(\omega)^6$ Klasseninvariante.}

In den Fällen, wo $m \equiv 3\ (\mathrm{mod}\ 4)$ ist, können wir noch einen Schritt weiter gehen. Wir haben in § 73 neben den \gesperrt{Schlaeflischen} Modulargleichungen auch die Gleichungen
\weq{19}{\Phi_n(u_1, v_1) = 0}
kennen gelernt, die zwischen
\weq{20}{u_1 = f_1(\omega)^3, \quad v_1 = \left(\frac{2}{a}\right) f_1\left(\frac{c+\partial\omega}{a}\right)^3}
bestehen, und aus § 73, (10), (15) ergibt sich, daß, wenn $n \equiv 7\ (\mathrm{mod}\ 8)$ vorausgesetzt wird, $\Phi_n(u_1, v_1)$ rational durch
\[
u_1 v_1, \quad u_1^8 + v_1^8
\]
dargestellt werden kann. Wenn wir also in (19)
\weq{21}{v_1 = f_1\left(\frac{c+\partial\omega}{a}\right)^3 = \pm f_2(\omega)^3, \quad u_1 = f_1(\omega)^3}
setzen, so wird
\[
u_1 v_1 = \frac{\pm \sqrt{2}^{3}}{f(\omega)^3},
\]
und $u_1^8 + v_1^8$ kann rational durch $f(\omega)^{24}$ ausgedrückt werden [§ 34, (11)], d. h. es geht
\[
\Phi_n(u_1, v_1) = 0
\]
in eine oder, nach der Wahl des Vorzeichens, zwei rationale Gleichungen für
\[
\xi = \sqrt{2}\, f(\omega)^3
\]
über, die wir mit
\weq{22}{F(\pm\xi) = 0}
bezeichnen wollen.
